\section{Experiments}
\label{sec:experiments}

\subsection{Experimental Setup}

\textbf{Dataset.} We collect 30 Chinese ink-wash (Shanshui)
paintings spanning the Song, Yuan, Ming, and Qing dynasty styles.
Each painting is annotated with: (1)~a line-distribution map
(structural draft), (2)~an ink-density distribution map (5-channel,
five-level ink), and (3)~a blank-space mask (negative-space
regions). Annotations are created semi-automatically: a U-Net
segmenter produces initial masks that are refined by a human
annotator.

\textbf{Baselines.} We compare four methods:
\begin{itemize}
    \item \textbf{Wan2.1-I2V (unconditional)}: the base I2V model
    with no art-specific conditioning.
    \item \textbf{Single-condition (line only)}: line-distribution
    map injected at all denoising steps.
    \item \textbf{Single-condition (ink only)}: ink-density map
    injected at all denoising steps.
    \item \textbf{InkFlow (full)}: our dual-conditioned method with
    blank-space loss and rhythm prior.
\end{itemize}

\textbf{Metrics.}
\begin{itemize}
    \item \textbf{Ink-density histogram matching} (lower = better):
    EMD between the generated frame's ink-level histogram and the
    target histogram from the ink-density map.
    \item \textbf{Blank-space texture ratio} (lower = better):
    fraction of blank-space pixels that contain generated texture
    above a variance threshold.
    \item \textbf{Temporal consistency} (higher = better): average
    optical-flow warping error between consecutive frames, measured
    via the warping error metric from \citet{lai2018video}.
    \item \textbf{FID} (lower = better): Fréchet Inception Distance
    between generated frames and real ink-wash paintings.
    \item \textbf{User study}: 20 participants rate each video on
    7-point Likert scales for ink fidelity, blank-space
    preservation, motion naturalness, and overall aesthetic quality.
\end{itemize}

\subsection{Quantitative Results}

\textbf{[Table placeholder: quantitative comparison of 4 methods ×
5 metrics on 30 paintings. To be filled upon experiment completion.]}

\subsection{Qualitative Results}

\textbf{[Figure placeholder: side-by-side video frames from 4 methods
on 3 representative paintings, showing ink-density fidelity and
blank-space preservation. To be filled upon experiment completion.]}

\subsection{Ablation Study}

\textbf{[Table placeholder: ablation removing each component
(dual-stage separation, blank-space loss, rhythm prior) and
measuring impact on metrics. To be filled upon experiment
completion.]}

\subsection{Perceptual User Study}

\textbf{[Results pending.]} 20 participants view videos from all
four methods in randomized order and rate them on 7-point Likert
scales. We analyze ratings via Mann--Whitney $U$ tests with
Bonferroni correction for multiple comparisons.

\subsection{Status of Pending Experiments}

The InkFlow framework is designed and the dual-condition adapter
architecture is implemented. The following steps are required to
complete the evaluation:
\begin{enumerate}
    \item Collect and annotate the 30-painting dataset (line maps,
    ink-density maps, blank-space masks).
    \item Fine-tune Wan2.1-I2V with the dual-condition adapters and
    art-specific losses on the OilPaintingVideo ComfyUI pipeline.
    \item Generate videos for all 30 paintings under all 4 methods.
    \item Compute quantitative metrics.
    \item Conduct the perceptual user study.
\end{enumerate}
We expect all experiments to be completed before the camera-ready
deadline.
